% TOP_PROBLEM: {"id": 2301041, "title": "Research Problems in Function Theory — Problem 1.41", "category_id": 8, "category": {"id": 8, "name": "analysis", "display_name": "Analysis", "description": "Limits, continuity, calculus, and function theory.", "slug": "analysis", "order_index": 8, "created_at": "2026-07-31T15:26:25.671Z"}, "status": "open", "problem_number": "AMR-022-1041", "set_id": 13}
% FIRST_POSTED: 2026-10-02
% ENTRY_KIND: statement_only
% UnsolvedMath ID 2301041; source catalogue code AMR-022-1041
% !TeX program = lualatex
% Definition and sources only; this file is not an LLM solution attempt.
\documentclass[11pt,a4paper]{article}
\usepackage[margin=25mm]{geometry}
\usepackage{fontspec}
\setmainfont{lmroman10-regular.otf}[BoldFont=lmroman10-bold.otf,
 ItalicFont=lmroman10-italic.otf,BoldItalicFont=lmroman10-bolditalic.otf]
\usepackage{amsmath,amssymb,mathtools}
\usepackage{unicode-math}
\setmathfont{latinmodern-math.otf}
\AtBeginDocument{\let\setminus\smallsetminus}
\usepackage{xurl,hyperref}
\hypersetup{colorlinks=true,urlcolor=blue}
\setcounter{secnumdepth}{0}
\setlength{\parindent}{0pt}
\setlength{\parskip}{5pt}
\setlength{\emergencystretch}{3em}
\begin{document}
\section*{Research Problems in Function Theory — Problem 1.41}\label{2301041}
{\small UnsolvedMath ID 2301041 · AMR-022-1041 · Analysis · AMR Open Problem Lists}
\subsection{Definitions and mathematical statement}
For a nonconstant meromorphic function $f$, let $n(t,a)$ count solutions of $f(z)=a$ in $|z|\le t$, with multiplicity; $a=\infty$ means poles. Put
\[
N(r,a)=\int_0^r\frac{n(t,a)-n(0,a)}t\,dt+n(0,a)\log r,
\quad T(r,f)=N(r,\infty)+\frac1{2\pi}\int_0^{2\pi}\log^+|f(re^{i\theta})|\,d\theta,
\]
where $\log^+x=\max(0,\log x)$. At a boundary radius $R\in\{1,\infty\}$ with $T(r,f)\to\infty$, define the deficiency
$\delta(a,f)=1-\limsup_{r\to R}N(r,a)/T(r,f)$. A deficient value has positive deficiency.
Let $(a_j)$ be a finite or countable sequence of distinct points of $\widehat{\mathbb C}$, and let prescribed numbers $d_j\in(0,1]$ satisfy
$\sum_jd_j\le2$. The disc inverse-deficiency problem asks for a meromorphic function $f$ on $\mathbb D$ such that
\[
\limsup_{r\uparrow1}\frac{T(r,f)}{\log(1/(1-r))}=\infty,
\quad \delta(a_j,f)=d_j\quad\text{for every }j,
\]
and $\delta(a,f)=0$ at every other value $a\in\widehat{\mathbb C}$. The infinity value is treated as the pole value, not as a limit of the finite prescribed points. One may equivalently prescribe only finite values and prescribe the deficiency at infinity separately.
Are the individual restrictions $d_j\le1$ and the total bound two sufficient for all such data? Exact values and absence of extra deficient values are part of the question; constructing lower bounds $\delta(a_j,f)\ge d_j$ alone does not answer it.
\subsection{Short English statement}
Can every admissible list of deficiencies be realized exactly by a sufficiently rapidly growing meromorphic function in the unit disc?
\subsection{Sources}
\begin{itemize}
\item[S1] ULAM.AI and the UnsolvedMath Contributors, supplied problems.json, record 2301041 (AMR-022-1041), AMR Open Problem Lists. Catalogue identity and source transcription; CC BY 4.0.
\url{https://huggingface.co/datasets/ulamai/UnsolvedMath}.
\item[S2] Hayman - Research Problems in Function Theory (2018), Problem 1.41. Original source indexed by AMR-022-1041.
\url{https://arxiv.org/abs/1809.07200}.
\end{itemize}
\end{document}
